\section{Limitations}
\label{sec:threats}

\paragraph{Client-Side vs.\ Server-Side Bidding}
Our causal oracle intercepts client-side \texttt{Prebid.js} auctions. While client-side bids reflect direct market demand, proprietary Server-to-Server (S2S) frameworks (e.g., Google Open Bidding, Amazon TAM) remain opaque behind publisher ad servers~\cite{pachilakis2019no, cook2020inferring}.

\paragraph{Usability Proxy: Exception Overhead}
We quantify usability cost via normalized runtime JavaScript exceptions ($\ucost$). While exceptions capture severed script dependencies and broken execution graphs~\cite{iqbal2020adgraph}, they are an indirect proxy for functional failure. 
We treat $\ucost$ as a reproducible initial baseline and encourage the community to propose alternative methodologies (e.g., visual layout shifts or synthetic user journey completion rates) to standardize web breakage benchmarks.

\paragraph{Synthetic Personas vs.\ Longitudinal Drift}
Our 14-seeder profiles initialize commercial DSP device graphs and trigger premium bid pacing~\cite{bashir2016tracing, oh2022cart}, but lack the longitudinal entropy of multi-month human browsing (e.g., latent logins and search histories)~\cite{zeber2020representativeness, lerner2016internet}.

\paragraph{Software Rasterization Trade-offs}
Mesa software rasterization (\texttt{swrast}) enforced determinism across cluster workers and eliminated GPU confounding, though commercial scripts may probe bare-metal GPU driver signatures~\cite{mowery2012pixel, cao2017cross} that software rasterization inherently normalizes.

\paragraph{Geographic Vantage}
Workloads operated from European academic IPs (Portugal); GDPR and ePrivacy enforcement may induce distinct bid density and consent gating compared to non-GDPR regions~\cite{papadogiannakis2021user, barysei2023investigating}.